%% max_print_line: default
% `tex.web` §1279: a `\message` joins the line already in progress.
%
%     if term_offset+length(s)>max_print_line-2 then print_ln
%     else if (term_offset>0)or(file_offset>0) then print_char(" ");
%     slow_print(s); update_terminal;
%
% So consecutive messages come out on one line with single spaces between
% them, and one that would overflow the margin starts a line of its own
% instead. We gave every message a line of its own, which made two of
% trip.log's diffs phantoms: a `{restoring …}` that TeX had put after a
% message on the same line looked misordered against ours on the next one,
% and chasing it cost a probe to find that our restore order was right all
% along.
%
% Only another message leaves a line open. Everything else this engine writes
% ends its line the way `print_nl` does — a `\tracingcommands` line closes with
% `end_diagnostic(false)`, an error with §90's two `print_ln`s — so a message
% after one of those starts fresh.
%
% This runs at TeX's own margin rather than the 1000 the other cases use,
% because where the break falls is the whole subject.
\catcode`\{=1 \catcode`\}=2
\tracingonline=1

% A report first, so what follows starts at column zero in both engines: TeX
% opens its log with a `(./file.tex` this engine does not write, and a message
% would otherwise join *that*.
\showthe\escapechar

% Three that fit: one line.
\message{alpha}\message{beta}\message{gamma}

% ...and one that does not: it starts its own, and the next joins *that*.
\message{The quick brown fox jumps over the lazy dog and keeps on running past the post}\message{tail}

% A report closes the line, so what follows begins a new one.
\showthe\escapechar
\message{after-a-report}
\message{joined}

% A report last, so TeX's closing `)` lands on a line of its own rather than on
% the end of `joined`.
\showthe\escapechar
\end
